\begin{proofbox}
	\textbf{\textcolor{CProof}{思路提示}}

	从椭圆定义出发，先把焦点弦 $PQ$ 分解为
	\[
		|PQ|=|PF_1|+|QF_1|,
	\]
	再把三边长度和整理成两个椭圆定义式的和。

	\medskip
	\textbf{\textcolor{CProof}{正式推导}}
	\begin{description}[style=nextline, leftmargin=2.8em, labelsep=0.5em, itemsep=0.55em]
		\item[\textbf{步骤一（椭圆定义）：}]
		      因为 $P,Q$ 均在椭圆上，所以
		      \[
			      |PF_1|+|PF_2|=2a,\qquad |QF_1|+|QF_2|=2a.
		      \]
		      \par\textbf{理由：} 椭圆上任意一点到两个焦点的距离之和恒为 $2a$。

		\item[\textbf{步骤二（弦长分解）：}]
		      因为 $F_1$ 是椭圆内部点，直线 $PQ$ 过 $F_1$，并与椭圆交于 $P,Q$ 两点，所以 $P,Q$ 分居 $F_1$ 两侧，即 $F_1$ 在线段 $PQ$ 上。
		      因此
		      \[
			      |PQ|=|PF_1|+|QF_1|.
		      \]
		      \par\textbf{理由：} 在线段上，整体长度等于两部分长度之和。

		\item[\textbf{步骤三（代入求和）：}]
		      \[
			      \begin{aligned}
				      C_{\triangle PF_2Q}
				       & = |PF_2|+|QF_2|+|PQ|              \\
				       & = |PF_2|+|QF_2|+|PF_1|+|QF_1|     \\
				       & = (|PF_1|+|PF_2|)+(|QF_1|+|QF_2|) \\
				       & = 2a+2a                           \\
				       & = 4a.
			      \end{aligned}
		      \]
		      \par\textbf{理由：} 将步骤一、步骤二的结果代入，并把属于同一点的两焦点距离配对。
	\end{description}

	\medskip
	\textbf{\textcolor{CProof}{退化情形说明}}

	若 $PQ$ 不是长轴弦，则 $P,F_2,Q$ 不共线，$C_{\triangle PF_2Q}$ 就是普通三角形周长。

	若 $PQ$ 恰为长轴弦，则 $P,F_2,Q$ 三点共线，此时不是非退化三角形；但按三边长度和
	\[
		|PF_2|+|QF_2|+|PQ|
	\]
	理解，上述推导仍然成立，结果仍为 $4a$。

	\medskip
	\textbf{\textcolor{CProof}{结论回扣}}

	因此，对于任意过焦点 $F_1$ 的弦 $PQ$，与另一焦点 $F_2$ 构成的三边长度和恒为 $4a$，与弦的位置无关。
\end{proofbox}